\chapter{کلیات پژوهش}

\section{مقدمه}
آشکارسازی شیء مسئله تعیین مکان و رده هر شیء در یک تصویر یا فریم ویدئویی است. خروجی آشکارساز معمولاً مجموعه‌ای از جعبه‌های مرزی، برچسب‌های رده و امتیازهای اطمینان است. این وظیفه با طبقه‌بندی تصویر تفاوت دارد؛ زیرا مدل باید تعداد متغیری از اشیا را در اندازه‌ها و مکان‌های گوناگون تشخیص دهد. وجود اشیای کوچک، هم‌پوشانی، انسداد، تغییر مقیاس و پس‌زمینه‌های شلوغ، آشکارسازی را به مسئله‌ای حساس و چندوجهی تبدیل می‌کند \cite{lin2014coco}.

روش‌های دو مرحله‌ای، مانند FasterR-CNN، ابتدا ناحیه‌های نامزد را تولید و سپس آن‌ها را طبقه‌بندی و مکان‌یابی می‌کنند \cite{ren2015faster}. در مقابل، آشکارسازهای تک‌مرحله‌ای مانند SSD و YOLO پیش‌بینی‌ها را مستقیماً از نقشه‌های ویژگی استخراج می‌کنند و برای کاربردهای بلادرنگ مناسب‌اند \cite{liu2016ssd,redmon2016yolo}. مسئله کلیدی در این خانواده، یافتن توازن میان ظرفیت بازنمایی و هزینه اجرا است: افزایش عمق، پهنای کانال یا پیچیدگی توجه ممکن است دقت را بالا ببرد، اما محدودیت تأخیر و حافظه را نقض کند.

YOLOv13 از کانولوشن‌های تفکیک‌پذیر برای کاهش هزینه محاسباتی بهره می‌گیرد و با HyperACE و FullPAD به تقویت هم‌بستگی‌های چندمقیاسی و جریان اطلاعات در کل مدل می‌پردازد \cite{lei2025yolov13}. در پژوهش حاضر، این خط پایه با دو جهت‌گیری تکمیل می‌شود. نخست، توجه Transformer-in-Transformer برای بازنمایی ساختارهای محلی درون قطعه و رابطه‌های میان قطعه بررسی می‌شود \cite{han2021tnt}. دوم، بلوک‌های Inception همراه با توجه کانالی کارا برای ترکیب پاسخ‌های چندمقیاسی و برجسته‌سازی کانال‌های اطلاع‌رسان به کار می‌روند \cite{szegedy2015inception,wang2020eca}.

\section{بیان مسئله}
مسئله اصلی این پژوهش، طراحی معماری تک‌مرحله‌ای برای آشکارسازی شیء است که دقت تشخیص را با هزینه محاسباتی و تأخیر مناسب همراه کند. سبک‌سازی بی‌هدف شبکه ممکن است اطلاعات موردنیاز برای تشخیص اشیای کوچک یا پوشیده را حذف کند. در سوی مقابل، افزودن گسترده بلوک‌های توجه و شاخه‌های چندمقیاسی می‌تواند تعداد پارامترها، FLOPs، مصرف حافظه و زمان استنتاج را افزایش دهد. بنابراین، انتخاب نوع ماژول، محل استقرار آن در شبکه، تعداد تکرار و شیوه ادغام ویژگی‌ها تعیین‌کننده است.

معماری مورد بررسی بر YOLOv13 استوار است و سه مؤلفه را به‌صورت هدفمند ترکیب می‌کند: توجه TNT برای تعاملات محلی-سراسری، بلوک اینسپشن مجهز به توجه کانالی برای پردازش چندمقیاسی، و توجه هندسی برای بررسی تابع شباهت در \lr{AAttn}. مسئله پژوهش آن است که کدام پیکربندی از این مؤلفه‌ها، بهترین توازن را میان کیفیت آشکارسازی و هزینه اجرایی ایجاد می‌کند.

\section{اهمیت و ضرورت پژوهش}
در کاربردهای عملی، مانند دوربین‌های هوشمند و سامانه‌های لبه، کیفیت آشکارسازی تنها معیار تصمیم نیست. مدلی که AP بالاتری دارد ولی به حافظه زیاد، توان محاسباتی بالا یا تأخیر غیرقابل قبول نیاز دارد، ممکن است برای محیط مقصد قابل استفاده نباشد. از این رو، تحلیل هم‌زمان AP، تعداد پارامتر، FLOPs، تأخیر، نرخ فریم و مصرف حافظه ضروری است. افزون بر این، اشیای کوچک و صحنه‌های متراکم به دریافت نشانه‌ها در میدان دیدهای متفاوت و مدل‌سازی ارتباط میان اجزای تصویر نیاز دارند. ترکیب توجه سلسله‌مراتبی با پردازش چندمقیاسی سبک‌وزن، پاسخی معماری به این نیاز است.

\section{اهداف پژوهش}
\subsection{هدف اصلی}
طراحی و ارزیابی آشکارساز شیء سبک و دقیق مبتنی بر YOLOv13 با استفاده هدفمند از توجه Transformer-in-Transformer و بلوک‌های Inception.

\subsection{اهداف فرعی}
\begin{itemize}
  \item بررسی اثر TNTBlock در بخش‌های مختلف ستون فقرات بر دقت و هزینه اجرا؛
  \item بررسی بلوک InceptionECA برای استخراج ویژگی چندمقیاسی و توجه کانالی؛
  \item مقایسه توجه هندسی با خط پایه توجه ضرب داخلی و نسخه‌های حذفی آن؛
  \item انتخاب پیکربندی برتر بر پایه AP، تعداد پارامتر، FLOPs، تأخیر و حافظه؛
  \item تدوین پروتکل آزمایشی بازتولیدپذیر برای مقایسه منصفانه با YOLOv13.
\end{itemize}

\section{پرسش‌های پژوهش}
\begin{enumerate}
  \item توجه Transformer-in-Transformer چه اثری بر کیفیت آشکارسازی و هزینه اجرا دارد؟
  \item بلوک InceptionECA در کدام بخش‌های شبکه بهترین نسبت دقت به هزینه محاسباتی را فراهم می‌کند؟
  \item کدام گونه توجه هندسی از نظر AP، پایداری آموزش و زمان استنتاج مناسب‌تر است؟
  \item معماری نهایی در مقایسه با YOLOv13 چه توازنی میان دقت، پارامتر، FLOPs، تأخیر و حافظه ایجاد می‌کند؟
\end{enumerate}

\section{روش پژوهش و طرح ارزیابی}
پژوهش حاضر از نوع کاربردی و مبتنی بر آزمایش‌های کنترل‌شده است. ابتدا YOLOv13 با تنظیمات ثابت آموزش و ارزیابی می‌شود. سپس اثر هر مؤلفه معماری به‌طور مستقل بررسی می‌گردد و در مرحله بعد ترکیب‌های منتخب مقایسه می‌شوند. مجموعه‌داده، تفکیک‌پذیری ورودی، تعداد دوره‌ها، افزایش‌داده، بهینه‌ساز، زمان‌بند نرخ یادگیری و سخت‌افزار در مقایسه‌های هر آزمایش ثابت نگه داشته می‌شوند.

معیار اصلی ارزیابی، AP میانگین در آستانه‌های IoU از ۰٫۵ تا ۰٫۹۵ است. AP50 و AP75 نیز برای تحلیل رفتار مدل در آستانه‌های سهل‌گیرانه و سخت‌گیرانه گزارش می‌شوند. هزینه مدل با تعداد پارامترها، FLOPs، تأخیر هر تصویر پس از گرم‌کردن، نرخ فریم و مصرف بیشینه حافظه اندازه‌گیری می‌شود. برای گزارش پایدار، نتایج اصلی در چند بذر تصادفی تکرار و میانگین و انحراف معیار آن‌ها ارائه خواهد شد.

\section{نوآوری پژوهش}
نوآوری پژوهش، طراحی و تحلیل یک آشکارساز سبک‌وزن است که سه ایده مکمل را یکپارچه می‌کند: توجه سلسله‌مراتبی درون‌قطعه‌ای و میان‌قطعه‌ای، استخراج ویژگی چندمقیاسی با کانولوشن‌های عمقی در بلوک‌های Inception، و تحلیل تابع شباهت هندسی در توجه. ارزش این طراحی با مطالعه حذفی شفاف و تحلیل مبادله دقت-هزینه سنجیده می‌شود.

\section{ساختار پایان‌نامه}
فصل اول مسئله، اهداف، پرسش‌ها و روش ارزیابی را معرفی کرد. فصل دوم مبانی و پیشینه آشکارسازی شیء، YOLO، مبدل‌ها و بلوک‌های چندمقیاسی را مرور می‌کند. فصل سوم معماری نهایی و جزئیات پیاده‌سازی را توضیح می‌دهد. فصل چهارم نتایج کمی، نتایج کیفی و مطالعات حذفی را ارائه می‌کند و فصل پنجم به جمع‌بندی و پیشنهادهای آینده می‌پردازد.
